% methodology.tex -- EVM Executor Paper (BNR-2026-008)

\section{Methodology}

\subsection{Research Design}

This study employs a three-phase methodology:

\begin{enumerate}
  \item \textbf{Comparative architecture analysis}: Systematic examination of
    five production L2 systems to extract common patterns and identify the
    minimal EVM fork surface.
  \item \textbf{Opcode constraint mapping}: Classification of all 75 explicitly
    analyzed Cancun EVM opcodes (plus ${\sim}$65 PUSH/DUP/SWAP variants, totaling
    ${\sim}$140) by ZK proving difficulty, estimated constraint count (R1CS and
    PLONKish), and recommended circuit treatment.
  \item \textbf{Trace format design and benchmark projection}: Design of a
    selective ZK tracer using Geth's \texttt{core/tracing} hooks, with throughput
    projections derived from published benchmarks.
\end{enumerate}

\subsection{Opcode Constraint Classification}

Each EVM opcode is classified along three dimensions:

\textbf{ZK Difficulty} (six tiers):
\begin{itemize}
  \item \textsc{Trivial} ($<$10 constraints): Pure stack manipulation (POP, PUSH,
    DUP, SWAP, PC).
  \item \textsc{Cheap} (10--100 constraints): Simple arithmetic (ADD, SUB, MUL),
    comparisons, memory loads.
  \item \textsc{Moderate} (100--1{,}000 constraints): Division, modular
    arithmetic, memory copies, log emissions.
  \item \textsc{Expensive} (1{,}000--10{,}000 constraints): Storage access
    (SLOAD, SSTORE), external code queries, balance lookups.
  \item \textsc{Very Expensive} (10{,}000--100{,}000 constraints): Cross-contract
    calls (CALL, DELEGATECALL, STATICCALL), contract creation (CREATE, CREATE2).
  \item \textsc{Prohibitive} ($>$100{,}000 constraints): KECCAK256.
\end{itemize}

\textbf{Circuit Treatment} (six strategies):
\begin{itemize}
  \item \textsc{Direct}: Prove directly in arithmetic circuit.
  \item \textsc{Lookup}: Use precomputed lookup tables.
  \item \textsc{Oracle}: Use preimage oracle for hash precomputation.
  \item \textsc{Special}: Dedicated sub-circuit required.
  \item \textsc{Replace}: Substitute with ZK-friendly alternative.
  \item \textsc{Unsupported}: Excluded from initial version.
\end{itemize}

\textbf{Constraint Count}: Estimated R1CS and PLONKish constraints per
invocation, sourced from Polygon zkEVM ZK counter
documentation~\cite{polygon2024zkcounters}, Scroll circuit
specifications~\cite{scroll2024whitepaper}, and the constraint-level
analysis of Burchardt et al.~\cite{burchardt2025constraint}.

\subsection{ZK Tracer Design}

The tracer implements Geth's \texttt{tracing.Hooks} struct with selective
callbacks:

\begin{lstlisting}[language=Go,caption={ZK Tracer hook registration}]
func (t *ZKTracer) Hooks() *tracing.Hooks {
  return &tracing.Hooks{
    OnOpcode:        t.onOpcode,
    OnStorageChange: t.onStorageChange,
    OnBalanceChange: t.onBalanceChange,
    OnNonceChange:   t.onNonceChange,
    OnLog:           t.onLog,
  }
}
\end{lstlisting}

In selective mode (default), the \texttt{OnOpcode} callback increments a
counter without recording individual instructions. In full-capture mode
(debugging only), it records the program counter, opcode name, remaining
gas, and call depth for each instruction.

The trace output structure captures:
\begin{itemize}
  \item \textbf{Storage accesses}: Address, slot, value for each SLOAD/SSTORE.
  \item \textbf{Balance changes}: Address, previous/new balance, change reason.
  \item \textbf{Nonce changes}: Address, previous/new nonce.
  \item \textbf{Log emissions}: Address, topics, data.
  \item \textbf{Opcode count}: Total instructions executed (always captured).
  \item \textbf{Contract creation}: Addresses of newly deployed contracts.
\end{itemize}

\subsection{Benchmark Projection Methodology}

Due to the Go runtime not being available on the experimental platform at the
time of this study, throughput benchmarks are projected from published data
using the following methodology:

\begin{enumerate}
  \item \textbf{Base throughput}: Derived from Geth's measured live sync rate of
    100--200~MGas/s~\cite{paradigm2024reth} and standard gas costs (21{,}000 gas
    for simple transfers, 20{,}000--22{,}000 gas for cold SSTORE operations).
  \item \textbf{In-memory adjustment}: A 1.5--2$\times$ multiplier accounts for
    in-memory StateDB eliminating disk I/O latency.
  \item \textbf{Tracing overhead}: Estimated from the difference between Geth's
    structLogger overhead (10--50$\times$) and the selective nature of our ZK
    tracer, yielding 10--30\% overhead for state-change-only capture.
  \item \textbf{Validation}: All projections are cross-referenced against
    independent benchmarks (Reth at 700~MGas/s single-threaded, Gravity Reth
    at 1.5~GGas/s, op-geth community benchmarks at 3{,}000--5{,}000~tx/s).
\end{enumerate}

Confidence levels are assigned to each projection: \textsc{High} (supported by
multiple independent sources), \textsc{Medium} (extrapolated from related
benchmarks), or \textsc{Low} (estimated from first principles with limited
validation data).
